\documentclass[12pt,a4paper]{article}

\usepackage[a4paper,text={16.5cm,25.2cm},centering]{geometry}
\usepackage{lmodern}
\usepackage{amssymb,amsmath}
\usepackage{bm}
\usepackage{graphicx}
\usepackage{microtype}
\usepackage{hyperref}
\usepackage[usenames,dvipsnames]{xcolor}
\setlength{\parindent}{0pt}
\setlength{\parskip}{1.2ex}




\hypersetup
       {   pdfauthor = {  },
           pdftitle={  },
           colorlinks=TRUE,
           linkcolor=black,
           citecolor=blue,
           urlcolor=blue
       }




\usepackage{upquote}
\usepackage{listings}
\usepackage{xcolor}
\lstset{
    basicstyle=\ttfamily\footnotesize,
    upquote=true,
    breaklines=true,
    breakindent=0pt,
    keepspaces=true,
    showspaces=false,
    columns=fullflexible,
    showtabs=false,
    showstringspaces=false,
    escapeinside={(*@}{@*)},
    extendedchars=true,
}
\newcommand{\HLJLt}[1]{#1}
\newcommand{\HLJLw}[1]{#1}
\newcommand{\HLJLe}[1]{#1}
\newcommand{\HLJLeB}[1]{#1}
\newcommand{\HLJLo}[1]{#1}
\newcommand{\HLJLk}[1]{\textcolor[RGB]{148,91,176}{\textbf{#1}}}
\newcommand{\HLJLkc}[1]{\textcolor[RGB]{59,151,46}{\textit{#1}}}
\newcommand{\HLJLkd}[1]{\textcolor[RGB]{214,102,97}{\textit{#1}}}
\newcommand{\HLJLkn}[1]{\textcolor[RGB]{148,91,176}{\textbf{#1}}}
\newcommand{\HLJLkp}[1]{\textcolor[RGB]{148,91,176}{\textbf{#1}}}
\newcommand{\HLJLkr}[1]{\textcolor[RGB]{148,91,176}{\textbf{#1}}}
\newcommand{\HLJLkt}[1]{\textcolor[RGB]{148,91,176}{\textbf{#1}}}
\newcommand{\HLJLn}[1]{#1}
\newcommand{\HLJLna}[1]{#1}
\newcommand{\HLJLnb}[1]{#1}
\newcommand{\HLJLnbp}[1]{#1}
\newcommand{\HLJLnc}[1]{#1}
\newcommand{\HLJLncB}[1]{#1}
\newcommand{\HLJLnd}[1]{\textcolor[RGB]{214,102,97}{#1}}
\newcommand{\HLJLne}[1]{#1}
\newcommand{\HLJLneB}[1]{#1}
\newcommand{\HLJLnf}[1]{\textcolor[RGB]{66,102,213}{#1}}
\newcommand{\HLJLnfm}[1]{\textcolor[RGB]{66,102,213}{#1}}
\newcommand{\HLJLnp}[1]{#1}
\newcommand{\HLJLnl}[1]{#1}
\newcommand{\HLJLnn}[1]{#1}
\newcommand{\HLJLno}[1]{#1}
\newcommand{\HLJLnt}[1]{#1}
\newcommand{\HLJLnv}[1]{#1}
\newcommand{\HLJLnvc}[1]{#1}
\newcommand{\HLJLnvg}[1]{#1}
\newcommand{\HLJLnvi}[1]{#1}
\newcommand{\HLJLnvm}[1]{#1}
\newcommand{\HLJLl}[1]{#1}
\newcommand{\HLJLld}[1]{\textcolor[RGB]{148,91,176}{\textit{#1}}}
\newcommand{\HLJLs}[1]{\textcolor[RGB]{201,61,57}{#1}}
\newcommand{\HLJLsa}[1]{\textcolor[RGB]{201,61,57}{#1}}
\newcommand{\HLJLsb}[1]{\textcolor[RGB]{201,61,57}{#1}}
\newcommand{\HLJLsc}[1]{\textcolor[RGB]{201,61,57}{#1}}
\newcommand{\HLJLsd}[1]{\textcolor[RGB]{201,61,57}{#1}}
\newcommand{\HLJLsdB}[1]{\textcolor[RGB]{201,61,57}{#1}}
\newcommand{\HLJLsdC}[1]{\textcolor[RGB]{201,61,57}{#1}}
\newcommand{\HLJLse}[1]{\textcolor[RGB]{59,151,46}{#1}}
\newcommand{\HLJLsh}[1]{\textcolor[RGB]{201,61,57}{#1}}
\newcommand{\HLJLsi}[1]{#1}
\newcommand{\HLJLso}[1]{\textcolor[RGB]{201,61,57}{#1}}
\newcommand{\HLJLsr}[1]{\textcolor[RGB]{201,61,57}{#1}}
\newcommand{\HLJLss}[1]{\textcolor[RGB]{201,61,57}{#1}}
\newcommand{\HLJLssB}[1]{\textcolor[RGB]{201,61,57}{#1}}
\newcommand{\HLJLnB}[1]{\textcolor[RGB]{59,151,46}{#1}}
\newcommand{\HLJLnbB}[1]{\textcolor[RGB]{59,151,46}{#1}}
\newcommand{\HLJLnfB}[1]{\textcolor[RGB]{59,151,46}{#1}}
\newcommand{\HLJLnh}[1]{\textcolor[RGB]{59,151,46}{#1}}
\newcommand{\HLJLni}[1]{\textcolor[RGB]{59,151,46}{#1}}
\newcommand{\HLJLnil}[1]{\textcolor[RGB]{59,151,46}{#1}}
\newcommand{\HLJLnoB}[1]{\textcolor[RGB]{59,151,46}{#1}}
\newcommand{\HLJLoB}[1]{\textcolor[RGB]{102,102,102}{\textbf{#1}}}
\newcommand{\HLJLow}[1]{\textcolor[RGB]{102,102,102}{\textbf{#1}}}
\newcommand{\HLJLp}[1]{#1}
\newcommand{\HLJLc}[1]{\textcolor[RGB]{153,153,119}{\textit{#1}}}
\newcommand{\HLJLch}[1]{\textcolor[RGB]{153,153,119}{\textit{#1}}}
\newcommand{\HLJLcm}[1]{\textcolor[RGB]{153,153,119}{\textit{#1}}}
\newcommand{\HLJLcp}[1]{\textcolor[RGB]{153,153,119}{\textit{#1}}}
\newcommand{\HLJLcpB}[1]{\textcolor[RGB]{153,153,119}{\textit{#1}}}
\newcommand{\HLJLcs}[1]{\textcolor[RGB]{153,153,119}{\textit{#1}}}
\newcommand{\HLJLcsB}[1]{\textcolor[RGB]{153,153,119}{\textit{#1}}}
\newcommand{\HLJLg}[1]{#1}
\newcommand{\HLJLgd}[1]{#1}
\newcommand{\HLJLge}[1]{#1}
\newcommand{\HLJLgeB}[1]{#1}
\newcommand{\HLJLgh}[1]{#1}
\newcommand{\HLJLgi}[1]{#1}
\newcommand{\HLJLgo}[1]{#1}
\newcommand{\HLJLgp}[1]{#1}
\newcommand{\HLJLgs}[1]{#1}
\newcommand{\HLJLgsB}[1]{#1}
\newcommand{\HLJLgt}[1]{#1}


\def\endash{–}
\def\bbD{ {\mathbb D} }
\def\bbZ{ {\mathbb Z} }
\def\bbR{ {\mathbb R} }
\def\bbC{ {\mathbb C} }

\def\x{ {\vc x} }
\def\a{ {\vc a} }
\def\b{ {\vc b} }
\def\c{ {\vc c} }
\def\e{ {\vc e} }
\def\k{ {\vc k} }
\def\f{ {\vc f} }
\def\g{ {\vc g} }
\def\h{ {\vc h} }
\def\u{ {\vc u} }
\def\v{ {\vc v} }
\def\y{ {\vc y} }
\def\z{ {\vc z} }
\def\w{ {\vc w} }

\def\bt{ {\tilde b} }
\def\ct{ {\tilde c} }
\def\Ut{ {\tilde U} }
\def\Qt{ {\tilde Q} }
\def\Rt{ {\tilde R} }
\def\Xt{ {\tilde X} }
\def\acos{ {\rm acos}\, }
\def\sinc{ {\rm sinc}\, }

\def\red#1{ {\color{red} #1} }
\def\blue#1{ {\color{blue} #1} }
\def\green#1{ {\color{ForestGreen} #1} }
\def\magenta#1{ {\color{magenta} #1} }

\input{somacros}

\begin{document}



\textbf{Numerical Analysis MATH50003 (2025\ensuremath{\endash}26) Revision Sheet}

\textbf{Problem 1} Show that
\[
{f(x+2h) - f(x-h) \over 3 h} = f'(x) + \ensuremath{\delta}
\]
where
\[
|\ensuremath{\delta}| \ensuremath{\leq} 5Mh/6
\]
for $M = \sup_{x-h \ensuremath{\leq} t \ensuremath{\leq} x+2h}|f''(t)|$

\textbf{Problem 2(a)} Consider a Lower triangular matrix with floating point entries:
\[
L = \begin{bmatrix}
\ensuremath{\ell}_{11} \\
 \ensuremath{\ell}_{21} & \ensuremath{\ell}_{22} \\
 \ensuremath{\vdots} & \ensuremath{\ddots} & \ensuremath{\ddots} \\
 \ensuremath{\ell}_{n1} & \ensuremath{\cdots} & \ensuremath{\ell}_{n,n-1} & \ensuremath{\ell}_{nn}
 \end{bmatrix} \ensuremath{\in} F_{\ensuremath{\sigma},Q,S}^{n \ensuremath{\times} n}
\]
and a vector $\ensuremath{\bm{\x}} \in F_{\ensuremath{\sigma},Q,S}^{n}$, where $F_{\ensuremath{\sigma},Q,S}$ is a set of floating-point numbers. Denoting matrix-vector multiplication implemented using floating point arithmetic as
\[
\ensuremath{\bm{\b}} := {\tt lowermul}(L,\ensuremath{\bm{\x}})
\]
express the entries $b_k := {\bf e}_k^\ensuremath{\top} \ensuremath{\bm{\b}}$  in terms of $\ensuremath{\ell}_{kj}$ and $x_k := {\bf e}_k^\ensuremath{\top} \ensuremath{\bm{\x}}$,  using rounded floating-point operations $\ensuremath{\oplus}$ and $\ensuremath{\otimes}$.

\textbf{Problem 2(b)} Assuming all operations involve normal floating numbers, show that your approximation has the form
\[
L \ensuremath{\bm{\x}} = {\tt lowermul}(L, \ensuremath{\bm{\x}}) + \ensuremath{\bm{\epsilon}}
\]
where, for $\ensuremath{\epsilon}_{\rm m}$ denoting machine epsilon and $E_{n,\ensuremath{\epsilon}}:= {n \ensuremath{\epsilon} \over 1-n\ensuremath{\epsilon}}$ and assuming $n \ensuremath{\epsilon}_{\rm m} < 2$,
\[
\| \ensuremath{\bm{\epsilon}} \|_1 \ensuremath{\leq}   2E_{n,\ensuremath{\epsilon}_{\rm m}/2}   \|L\|_1 \| \ensuremath{\bm{\x}} \|_1.
\]
Here we use  the matrix norm $\| A \|_1 := \max_j \ensuremath{\sum}_{k=1}^n |a_{kj}|$ and the vector norm $\| \ensuremath{\bm{\x}} \|_1 := \ensuremath{\sum}_{k=1}^n |x_k|$. You may use the fact that
\[
x_1 \ensuremath{\oplus} \ensuremath{\cdots} \ensuremath{\oplus} x_n = x_1 +  \ensuremath{\cdots} + x_n + \ensuremath{\sigma}_n
\]
where
\[
|\ensuremath{\sigma}_n| \ensuremath{\leq} \| \ensuremath{\bm{\x}} \|_1 E_{n-1,\ensuremath{\epsilon}_{\rm m}/2}.
\]
\textbf{Problem 3} What is the dual extension of square-roots? I.e. what should $\sqrt{a + b \ensuremath{\epsilon}}$ equal assuming $a > 0$?

\textbf{Problem 4} Use the Cholesky factorisation to determine whether the following matrix is symmetric positive definite:
\[
\begin{bmatrix} 2 & 2 & 1  \\
2 & 3 & 2\\
1 & 2 & 2
\end{bmatrix}
\]
\textbf{Problem 5} Compute the QR factorisation of
\[
\begin{bmatrix}
-3 & 1 \\
4 & 2
\end{bmatrix}
\]
\textbf{Problem 6} For the function $f(\ensuremath{\theta}) = \sin 3 \ensuremath{\theta}$, state explicit formulae for its Fourier coefficients
\[
\hat f_k := {1 \over 2\ensuremath{\pi}} \int_0^{2\ensuremath{\pi}} f(\ensuremath{\theta}) {\rm e}^{-{\rm i} k \ensuremath{\theta}} {\rm d}\ensuremath{\theta}
\]
and  their discrete approximation:
\[
\hat f_k^n := {1 \over n} \sum_{j=0}^{n-1} f(\ensuremath{\theta}_j) {\rm e}^{-{\rm i} k \ensuremath{\theta}_j}.
\]
for \emph{all} integers $k$, $n = 1,2,\ensuremath{\ldots}$, where $\ensuremath{\theta}_j = 2\ensuremath{\pi} j/n$.

\textbf{Problem 7(a)} Prove that the second kind Chebyshev polynomials
\[
U_n(x) := {\sin (n+1) \acos x \over \sqrt{1-x^2}}
\]
satisfy for any polynomial $p(x)$ of degree $m < n$
\[
\int_{-1}^1 U_n(x) p(x) \sqrt{1-x^2} {\rm d}x.
\]
\textbf{Problem 7(b)} Use this to prove that
\[
C_n(x) := U_{n+1}'(x)
\]
are orthogonal with respect to the inner product
\[
\ensuremath{\int}_{-1}^1 f(x)g(x) (1-x^2)^{3/2} {\rm d}x.
\]


\end{document}